\section{A variational principle for two-call reset experiments}
\label{sec:resetvariational91}
The fresh-preparation cut admits an exact finite-dimensional description
of optimal decision values.  No symmetry, second-moment identity, or
special choice of prior is needed in this section.  The description is
by a concave envelope over density matrices, not over pure states alone.
Its finite realization bounds the receiver needed to attain a value;
it does not identify reset width with a prescribed memory dimension.

\subsection{Experiments and strategy classes}
Let $I_\ell=\C^{d_\ell}$ and let $Y_\ell$ be a finite classical alphabet,
$\ell=1,2$.  A finite experiment consists of a probability law
$(p_\theta)$, two ordered measurements $E^{\theta,\ell}$ on $I_\ell$,
and rewards $r_{j\theta}\in[0,1]$ for decisions $1\le j\le m$.
The index $\theta$ is sampled once.  Conditional on it the two slots are
the memoryless measurement channels $\mathcal M_{E^{\theta,\ell}}$.
Repeating one unknown measurement is the case
$E^{\theta,1}=E^{\theta,2}$.  Identification and binary decisions on
subfamilies are obtained by taking indicator rewards.

\begin{definition}[Two-call classes]\label{def:classes91}
An unrestricted strategy prepares a state on $I_1\otimes R$, observes
the first label $y$, applies a common trace-preserving map from $R$ to
$I_2\otimes R'$, and, after observing $z$, measures $R'$ to decide $j$.
A unit-reset strategy instead applies a normal instrument
$(\mathcal K_{yh})_h$ to $R$, retains its old quantum output $O_{yh}$,
and prepares a state $\sigma_{yh}$ on the complete fresh system
$I_2\otimes R_{2,yh}$.  Conditional on $(y,h)$ the preparation is
exactly the tensor product of the old output with $\sigma_{yh}$.
No old quantum register enters $I_2$ or its fresh reference.  After the
second label, the two receivers may be measured jointly.  A classically
adaptive strategy is a unit-reset strategy that measures the entire old
receiver and retains only a classical record before fresh preparation.

All rules are independent of $\theta$ and defined on every formal
history, including histories null under a hypothesis.  Receiver spaces
have no dimension bound.  Instruments may have countably many outcomes;
all classical records are retained.  Public randomization and the
finite-dimensional trace-norm closure of the resulting output-dephased
Choi tester tuples are included in each class.  Early stopping is
implemented by a fixed, discarded fresh call and a decision independent
of its label; it is not a zero-length component of a prescribed profile.
Write $P_{\rm all},P_{\rm reset},P_{\rm ca}$ for the suprema of expected
reward in these classes.
\end{definition}
This reset cut is the two-call instance of Definition~\ref{def:reset88}.
The first class in the ordering
$P_{\rm ca}\le P_{\rm reset}\le P_{\rm all}$ agrees with complete
measure-and-reprepare adaptation, not with a constraint of separability
on all final reset effects.  In the order $I_1,Y_1,I_2,Y_2$, a
classically adaptive effect is a sum of positive products across
$I_1Y_1\mid I_2Y_2$ \cite[Definition~23]{OZNPQ89}.
The positive separable cone is closed in these finite dimensions:
its trace-one slice is the compact convex hull of product states, and
trace controls passage to the cone.  Only this necessary condition,
not its sufficiency for classical adaptive normalization, is used below.

\begin{lemma}[Countable records]\label{lem:countable91}
Let $(X_h)$ be positive trace-class operators with
$\sum_h\tr X_h<\infty$.  Their partial sums converge in trace norm,
and their recorded direct sum lies in the trace-class $\ell^1$ sum.
For effects $0\preceq F_h\preceq I$, the corresponding probability
series converges absolutely.  Replacing all but finitely many branches
of a normal instrument by one remainder branch preserves normalization;
changing the continuation on that branch changes any reward in $[0,1]$
by at most its probability.
\end{lemma}
\begin{proof}
The norm of a positive tail equals its trace.  This proves the first
assertions and bounds the probability tail by $\sum_{h>n}\tr X_h$.
For an instrument, the sum of the omitted completely positive maps is
the remainder map; its trace is precisely the omitted outcome
probability.  Apply the bound to each member of the finite experiment
and then average.  The maximum of its finitely many tail probabilities
tends to zero.  No conditional state is divided by a history probability.
\end{proof}

\subsection{A complete branch normal form}
Fix bases on $I_1,I_2$.  A coefficient map $C:I_1\to R$ represents
$|C\rangle=\sum_i|i\rangle\otimes C|i\rangle$; its squared norm is
$\tr C^*C$.  Measuring its input gives the unnormalized receiver state
$CE_y^{\mathsf T}C^*$.

\begin{lemma}[Unit-reset normal form and converse]\label{lem:resetnormal91}
Every unit-reset strategy has a purification, still respecting the reset
cut, from which its original decision statistics are recoverable by
common receiver channels.  For this purification there are rectangular
maps
\[
 C_0:I_1\longrightarrow R_0,\quad
 C_{yh}:I_1\longrightarrow O_{yh},\quad
 D_{yh}:I_2\longrightarrow R_{2,yh}
\]
with, on the input copies,
\begin{equation}\label{eq:resetgrams91}
 \rho=C_0^*C_0,\quad \tr\rho=1,\quad
 A_{yh}=C_{yh}^*C_{yh},\quad \sum_hA_{yh}=\rho,\quad
 b_{yh}=D_{yh}^*D_{yh},\quad\tr b_{yh}=1.
\end{equation}
The receiver block on the formal record $(y,h,z)$ is
\begin{equation}\label{eq:resetblock91}
 (C_{yh}E_y^{\theta,1\,\mathsf T}C_{yh}^*)
 \otimes(D_{yh}E_z^{\theta,2\,\mathsf T}D_{yh}^*).
\end{equation}
The sum in \eqref{eq:resetgrams91} converges in trace norm.
Conversely, every finite collection satisfying \eqref{eq:resetgrams91}
is realizable, up to common receiver isometries, by a unit-reset strategy
with old receiver dimension at most $d_1$ and fresh reference dimension
at most $d_2$.  Zero-probability histories require no extra assumption.
\end{lemma}
\begin{proof}
Purify the first preparation, retaining the purifying system in $R_0$.
The Schmidt support has dimension at most $d_1$.  For each $y$, dilate
each branch map of the intervening instrument to a single map
$V_{yh}:R_0\to O_{yh}$, retaining its environment as old receiver
memory.  On the relevant support,
$\sum_hV_{yh}^*V_{yh}=I$ in the strong sense.  Thus
$C_{yh}=V_{yh}C_0$ has
$\sum_h C_{yh}^*C_{yh}=C_0^*C_0$.
The positive tails converge in trace norm by Lemma~\ref{lem:countable91}.
This argument uses only the finite Schmidt support even when the
original receiver or a dilation environment is infinite-dimensional.
Tracing those environments recovers the original instrument.

Purify the entire fresh preparation separately.  Its coefficient map
is $D_{yh}$ with squared Hilbert--Schmidt norm one.  The tensor-product
preparation rule, not a statement about a probe marginal, gives
\eqref{eq:resetblock91}.  Common pre-query controls are absorbed into
these preparations and common post-query controls into the final
receiver measurement.  Fresh mixing may be purified inside the fresh
reference.  A retained public seed may be put in the classical record;
randomization before the first call is a mixture of such strategies and
cannot increase the supremum of a linear reward.  Countable records
and tester limits follow from the preceding lemma and continuity.

For the converse let $r=\operatorname{rank}\rho$ and choose
$C_0:I_1\to\C^r$ with $C_0^*C_0=\rho$.  Its restriction to
$\operatorname{supp}\rho$ is invertible onto $\C^r$.  Let $C_0^+$
be the inverse on this support, and set
\[
 C_{yh}=\sqrt{A_{yh}},\qquad
 V_{yh}=\sqrt{A_{yh}}C_0^+,\qquad D_{yh}=\sqrt{b_{yh}}.
\]
Positivity and $\sum_hA_{yh}=\rho$ imply
$\operatorname{supp}A_{yh}\subseteq\operatorname{supp}\rho$.
Hence $V_{yh}C_0=\sqrt{A_{yh}}$ and
$\sum_hV_{yh}^*V_{yh}=I_r$.  These are Kraus operators of a
trace-preserving recorded instrument with outputs in $\C^{d_1}$.
The fresh vector $|\sqrt{b_{yh}}\rangle$ is normalized and independent
of that output.  Polar decomposition identifies any other rectangular
maps of the same Gram matrices with these maps on their supports;
receiver isometries may be absorbed into the final effects.
Zero Gram matrices give zero branch maps.  All other null histories are
handled by the same unnormalized formulas, including projective devices.
\end{proof}

\subsection{The compact primal problem}
Let $\mathcal D_d$ denote the density matrices on $\C^d$, and put
\begin{equation}\label{eq:gameblocks91}
 G_{j,yz}=\sum_\theta p_\theta r_{j\theta}
     (E_y^{\theta,1})^{\mathsf T}\otimes(E_z^{\theta,2})^{\mathsf T}.
\end{equation}
These are fixed positive matrices on the two input copies.  For
$A\succeq0$ and $b\in\mathcal D_{d_2}$ define
\begin{equation}\label{eq:leafvalue91}
 \Phi_y(A,b)=\sum_z\max_{\substack{F_j\succeq0\\\sum_jF_j=I}}
 \sum_j\tr\bigl[F_j(\sqrt A\otimes\sqrt b)
            G_{j,yz}(\sqrt A\otimes\sqrt b)\bigr].
\end{equation}
The POVM is chosen separately at every recorded $z$.  It acts on
$\C^{d_1}\otimes\C^{d_2}$.  For $a\in\mathcal D_{d_1}$ let
$g_y(a)=\max_{b\in\mathcal D_{d_2}}\Phi_y(a,b)$ and let
\begin{equation}\label{eq:roof91}
 c_y(\rho)=\sup\left\{\sum_h\lambda_h g_y(a_h):
 \lambda_h\ge0,\ \sum_h\lambda_h=1,\
 \sum_h\lambda_ha_h=\rho\right\}.
\end{equation}
The envelope is taken over \emph{all} density matrices $a_h$.  Replacing
them by rank-one states would change the resource problem.  The
construction of concave envelopes is standard \cite{Uhlmann91}; the
operational content here is the reset realization and the common
barycenter constraint across all first labels.

\begin{theorem}[Exact reset variational principle]\label{thm:resetvariational91}
For every finite experiment of Definition~\ref{def:classes91},
\begin{equation}\label{eq:resetvariational91}
 P_{\rm reset}=\max_{\rho\in\mathcal D_{d_1}}\sum_yc_y(\rho)
 =\max_{\substack{A_{yh}\succeq0,\ \sum_hA_{yh}=\rho\\
                    \rho\in\mathcal D_{d_1},\ b_{yh}\in\mathcal D_{d_2}}}
                    \sum_{y,h}\Phi_y(A_{yh},b_{yh}).
\end{equation}
All maxima are attained.  For a fixed $\rho$ of rank $r$, at most $r^2$
nonzero instrument outcomes are required after each first label.
An optimal strategy exists with old receiver dimension at most $d_1$,
fresh reference dimension at most $d_2$, and a final quantum receiver
of dimension at most $d_1d_2$, separately from the finite classical
record.  These assertions include arbitrary priors and arbitrary fixed
pairs as particular decision problems; they assert an exact optimum,
not a strict advantage for every such problem.
\end{theorem}
\begin{proof}
For fixed $y$, $\Phi_y$ is continuous, and is positively homogeneous
in $A$.  Continuity follows from continuity of matrix square roots and
maximization of a continuous function over a compact POVM set.
Maximizing also over compact $\mathcal D_{d_2}$ makes $g_y$ continuous.
The graph $\{(a,g_y(a)):a\in\mathcal D_{d_1}\}$ is compact in a
real space of affine dimension at most $d_1^2$.  Its convex hull is
compact.  The largest height above $\rho$ is therefore attained and
is exactly $c_y(\rho)$, also allowing countable decompositions.
It is an upper semicontinuous concave function of $\rho$.
Consequently $\sum_yc_y$ attains a maximum.

Carath\'eodory first supplies at most $r^2+1$ atoms at a maximizing
height over a rank-$r$ matrix: each atom is supported on
$\operatorname{supp}\rho$, since positive matrices cannot cancel on its
orthogonal complement.  If more than $r^2$ weights are nonzero, the
Hermitian atoms are linearly dependent.  Choose nonzero real $v_h$ with
$\sum_hv_ha_h=0$; tracing gives $\sum_hv_h=0$.
Small perturbations $\lambda_h\pm\epsilon v_h$ preserve the barycenter
and positivity.  Optimality forces $\sum_hv_hg_y(a_h)=0$, since otherwise
one sign raises the height.  Move in this direction until a weight
vanishes.  The value is unchanged.  Iteration gives at most $r^2$ atoms.
For each atom the maximum over $b$ and the leaf POVMs is attained.

By Lemma~\ref{lem:resetnormal91}, the value of any purified competitor
is bounded by the right side of \eqref{eq:resetvariational91}, with
$A_{yh}=\lambda_{yh}a_{yh}$.  Countable records follow by the positive
tail bound; public mixtures and norm limits preserve the upper.
Conversely, the lemma realizes every finite maximizing collection,
and its optimal leaf POVMs realize every maximum in
\eqref{eq:leafvalue91}.  This proves equality and the dimension bounds.
The final quantum dimension excludes the separately recorded labels.
\end{proof}
For a binary equal-prior fixed pair, its optimal unhalved final trace
separation is $4P_{\rm reset}-2$.  Thus the theorem is also a general
variational formula for that operational distance.  The primal has
finitely many bounded matrix variables, but is not asserted to be a
convex semidefinite program or an efficient synthesis algorithm.

\subsection{Dual certificates and the role of receiver feedback}
\begin{theorem}[Duality and contact conditions]\label{thm:resetdual91}
For the same experiment,
\begin{equation}\label{eq:resetdual91}
 P_{\rm reset}=\min_{(H_y)}\lambda_{\max}\left(\sum_yH_y\right),
 \qquad \tr(H_ya)\ge g_y(a)\quad(a\in\mathcal D_{d_1}),
\end{equation}
where $H_y$ are Hermitian $d_1\times d_1$ matrices.  The minimum is
attained.  It is equivalent to minimizing $\lambda$ subject to the
same majorizations and $\sum_yH_y=\lambda I$.
For any optimal primal and dual, every atom of positive weight is a
contact point, $\tr(H_ya_{yh})=g_y(a_{yh})$, and $\rho$ is supported
on the maximal eigenspace of $\sum_yH_y$.  Conversely these conditions,
with maximizing fresh states and leaf POVMs, certify optimality.
\end{theorem}
\begin{proof}
Every majorization of $g_y$ majorizes its concave envelope.  Therefore
$\sum_yc_y(\rho)\le\tr(\rho\sum_yH_y)\le
\lambda_{\max}(\sum_yH_y)$, which proves weak duality.
For completeness, the strong-duality step is finite-dimensional.
Extend $c_y$ homogeneously to positive matrices by
$\widehat c_y(A)=\tr A\,c_y(A/\tr A)$ and $\widehat c_y(0)=0$.
Its hypograph
\[
 \mathcal K_y=\{(A,u):A\succeq0,\ u\le\widehat c_y(A)\}
\]
is a closed convex cone.  Closedness at $A=0$ follows from boundedness
of $g_y$; elsewhere it follows from the compact-graph construction.
The conic problem maximizing $\sum_yu_y$, with $(A_y,u_y)\in\mathcal K_y$,
$A_y=\rho$ and $\tr\rho=1$, is precisely the primal problem.
The point $A_y=\rho=I/d_1$, $u_y=-1$ is strictly interior to all
hypographs, since rewards are nonnegative and $A_y$ is positive definite.
Its value is finite, in fact at most one by the operational realization.
The finite-dimensional separating-hyperplane form of conic Slater
duality thus gives an attained dual with no gap
\cite[Section~5.2.3]{BV91}.

Its multipliers can be read directly: the supremum of
$u-\tr(H_yA)$ on $\mathcal K_y$ is zero exactly when
$\tr(H_ya)\ge g_y(a)$ for every density matrix.  Taking $\rho$ as a
free Hermitian variable under its trace constraint imposes
$\sum_yH_y=\lambda I$.  This gives the stated dual.  Conversely, for
majorants with sum $S$, add the positive matrix
$\lambda_{\max}(S)I-S$ to one $H_y$; all majorizations remain valid
and the sum becomes scalar.  The two forms have the same minimum.
Finally, the weak-duality difference is a sum of nonnegative atom
majorization defects and a nonnegative spectral defect.  It vanishes
exactly under the stated contact and support conditions.
\end{proof}
The majorizations in \eqref{eq:resetdual91} are universal inequalities
on density matrices.  A finite sampling of them is not a dual certificate.

\begin{corollary}[Label feedback versus receiver feedback]\label{cor:feedbackroof91}
If no receiver-instrument outcome is used to choose the second
preparation, but the first device label may be used, the optimum is
\[
 P_{\rm label}=\max_{\rho\in\mathcal D_{d_1}}\sum_yg_y(\rho).
\]
Allowing receiver feedback replaces $g_y$ by $c_y$ in this formula.
In particular, concavity of every $g_y$ is sufficient for
$P_{\rm label}=P_{\rm reset}$.  More generally, a label-only strategy
with a common density $\rho$ is optimal for the full reset class if
it attains the contact and spectral conditions of
Theorem~\ref{thm:resetdual91}.
\end{corollary}
\begin{proof}
With one preparation choice after each $y$, the normal form has
$A_y=\rho$; a trace-preserving receiver operation cannot improve a
subsequent unrestricted joint measurement and can be postponed to it.
Conversely these choices are admissible.  The assertion follows from
the primal formula and the definition of concave envelope.
\end{proof}

\subsection{The unrestricted causal normalization}
\begin{proposition}[Physical two-call tester blocks]\label{prop:causalnormal91}
In the unnormalized convention
$J^\sharp(E)=\sum_yE_y\otimes|y\rangle\langle y|$, every unrestricted
strategy has positive blocks $T_{j,yz}$ and positive matrices $\Xi_y$ with
\begin{equation}\label{eq:causalnormal91}
 \sum_jT_{j,yz}=\Xi_y,\qquad
 \tr_{I_2}\Xi_y=\rho_1,\qquad\tr\rho_1=1.
\end{equation}
Its probability of decision $j$ on the output block $(y,z)$ is
$\tr[T_{j,yz}(E_y^{\theta,1}\otimes E_z^{\theta,2})]$.
In particular $0\preceq T_{j,yz}\preceq\Xi_y$ and $\tr\Xi_y=1$.
\end{proposition}
\begin{proof}
Purify the first preparation to $|\psi\rangle\in I_1\otimes R$.
For its common intervening channel $\mathcal A_y:R\to I_2\otimes R'$
put
\[
 X_y=(\operatorname{Id}\otimes\mathcal A_y)(|\psi\rangle\langle\psi|),
 \quad\Xi_y=\tr_{R'}X_y,\quad\rho_1=\tr_R|\psi\rangle\langle\psi|.
\]
For the final POVM $(F_{j,yz})_j$ on $R'$ define
\[
 T_{j,yz}=\tr_{R'}[(I\otimes\sqrt{F_{j,yz}})X_y
                              (I\otimes\sqrt{F_{j,yz}})].
\]
Positivity is immediate.  Cyclicity in the traced factor and
$\sum_jF_{j,yz}=I$ give the sum; trace preservation of $\mathcal A_y$
gives the marginal.  Expanding the Born rule gives the contraction.
Here $\rho_1$ is the physical input marginal, whereas the coefficient
Gram matrix $\rho$ in \eqref{eq:resetgrams91} is its transpose.
These two conventions must not be interchanged.  Every complementary
effect is positive by the same construction.  No factor $d_1d_2$
occurs, because $J^\sharp$ uses the unnormalized maximally entangled
vector.  Countable records and limits follow by continuity.
\end{proof}
